\documentclass[dvipdfmx,11pt]{article}
\usepackage[top=28mm,bottom=28mm,left=34mm,right=34mm]{geometry}
%\renewcommand\baselinestretch{0.95}
\usepackage{amsmath,amsthm,amssymb,mathrsfs}
\usepackage{newtxtext,newtxmath}
\usepackage{enumitem}
\usepackage{tabularx}
%\usepackage{ascmac,braket,bm,pifont}
%\usepackage{fouridx}
%\pagestyle{empty}

\newcommand{\dayhead}[1]{\begin{center}\textsc{#1}\end{center}}
\newcommand{\talk}[4]{\begin{flushleft}#1 \textbf{#2} (#3)\\ \textit{#4}\end{flushleft}}
\newcommand{\abshead}[3]{\begin{flushleft}\textbf{#1} (#2)\\ \textit{#3}\end{flushleft}}
\newenvironment{abspar}{}{\vspace{.5\baselineskip}}

\begin{document}

\begin{center}
\Large
Hayama Symposium on\\
Complex Analysis in Several Variables XX\\
\& \\
Pacific Rim Complex-Symplectic Geometry Conference
\end{center}

\begin{center}
\vspace{\baselineskip}
\Large
\textsc{Program}
\end{center}

\vspace{\baselineskip}

\dayhead{Friday, July 13}

\talk{15:00--16:00}{Franc Forstneri\v{c}}{Ljubljana}{Runge Tubes}

\talk{16:30-17:30}{Yuta Kusakabe}{Osaka University}{Holomorphic flexibility properties of mapping spaces}

\vspace{\baselineskip}

\dayhead{Saturday, July 14}

\talk{9:30--10:30}{Youngjin Bae}{Kyoto University}{A Chekanov-Eliashberg algebra for Legendrian graphs}

\talk{10:45--11:45}{Tatsuki Kuwagaki}{University of Tokyo}{Irregular Riemann-Hilbert correspondence, irregular perverse sheaves, and Fukaya category}

\begin{flushleft}
Lunch
\end{flushleft}

\talk{14:00--15:00}{Chengjian Yao}{Universite libre de Bruxelles}{A cohomogeneity-one hypersymplectic flow and some convergence results}

\talk{15:30--16:30}{Eiji Inoue}{University of Tokyo}{Moduli space of Fano manifolds admitting K\"ahler-Ricci solitons}

\talk{16:45--17:45}{Ngoc Cuong Nguyen}{POSTECH}{H\"older continuous subsolution theorem for the complex Monge-Amp\`ere equation}

\clearpage

\dayhead{Sunday, July 15}

\talk{9:30--10:30}{Tien-Cuong Dinh}{National University of Singapore}{Sullivan's dictionary and ergodic theory for foliations}

\talk{10:45--11:45}{Sean Curry}{UC San Diego}{Strictly pseudoconvex domains in $\mathbb{C}^2$ with obstruction flat boundary}

\begin{flushleft}
Lunch
\end{flushleft}

\talk{14:00--15:00}{Junwu Tu}{University of Missouri}{Gromov-Witten invariants of matrix factorizations}

\talk{15:30--16:30}{Morimichi Kawasaki}{Kyoto University}{Poisson bracket invariant and the Lagrangian Floer theory}

\begin{flushleft}
Party
\end{flushleft}

\vspace{\baselineskip}

\dayhead{Monday, July 16}

\talk{9:30--10:30}{Yunhyung Cho}{Sungkyunkwan University}{Monotone Lagrangian tori in cotangent bundles of spheres}

\talk{10:45--11:45}{Kei Irie}{University of Tokyo}{Chain level loop bracket and pseudo-holomorphic curves}

\begin{flushleft}
Lunch
\end{flushleft}

\talk{14:00--15:00}{Damian Brotbek}{University of Strasbourg}{Hyperbolicity properties of general complete intersections}

\talk{15:30--16:30}{Tatsuo Suwa}{Hokkaido University}{Representation of Sato hyperfunctions by relative Dolbeault classes}

\talk{16:45--17:45}{Kyoung-Seog Lee}{IBS-CGP}{Special curves on some algebraic surfaces}

\vspace{\baselineskip}

\dayhead{Tuesday, July 17}

\talk{9:30--10:30}{Hiroshi Iriyeh}{Ibaraki University}{Mahler's conjecture for three dimensional convex bodies and the rigidity of certain symplectic embeddings}

\talk{10:45--11:45}{Filippo Bracci}{University of Roma Tor Vergata}{Continuous extension of biholomorphisms in higher dimension}


\clearpage

\begin{center}
\Large
\textsc{Abstracts}
\end{center}

\vspace{\baselineskip}

\dayhead{Friday, July 13}

\abshead{Franc Forstneri\v{c}}{Ljubljana}{Runge Tubes}

\begin{abspar}
I shall present a simple proof of the existence and plenitude of Runge 
tubes in complex Euclidean spaces $\mathbb{C}^n$ for $n>1$ and, more generally, in 
Stein manifolds with Varolin's density property. In particular, the 
inclusion of any closed complex submanifold $X$ into $\mathbb{C}^n$ can be 
approximated uniformly on compacts in $X$ by Runge embeddings of its 
normal bundle into $\mathbb{C}^n$.

\end{abspar}

\abshead{Yuta Kusakabe}{Osaka University}{Holomorphic flexibility properties of mapping spaces}

\begin{abspar}
Oka manifolds have holomorphic flexibility properties of $\mathbb{C}$-connectedness and dominability.
We generalize these properties to spaces of holomorphic maps and characterize Oka manifolds by holomorphic flexibility properties of their mapping spaces.
As an application, we shall prove a localization principle for the Oka property which gives new examples of Oka manifolds.

\end{abspar}

\vspace{\baselineskip}

\dayhead{Saturday, July 14}

\abshead{Youngjin Bae}{Kyoto University}{A Chekanov-Eliashberg algebra for Legendrian graphs}
\begin{abspar}
We define a differential graded algebra for Legendrian graphs in the standard contact Euclidean three space. This invariant is defined combinatorially by using ideas from the bordered version of Legendrian contact homology. A set of countably many generators and a generalized notion of equivalence are introduced for invariance. I will also talk about the relation between the DGA for Legendrian graph and a partially wrapped Floer homology of an induced Liouville sector. This is a joint work with Byung Hee An.
\end{abspar}

\abshead{Tatsuki Kuwagaki}{University of Tokyo}{Irregular Riemann-Hilbert correspondence, irregular perverse sheaves, and Fukaya category}

\begin{abspar}
Recently, D'Agnolo-Kashiwara established Riemann-Hilbert correspondence for possibly irregular holonomic D-modules, using R-constructible enhanced ind-sheaves. To see C-constructibility explicitly, we reformulate their equivalence using modules over "finite Novikov ring". As a byproduct, we can define irregular perverse sheaves generalizing usual perverse sheaves in a straightforward way. We also discuss a conjectural further reformulation of RH-correspondence in terms of Fukaya category which will explain the appearance of Novikov ring.
\end{abspar}

\clearpage

\abshead{Chengjian Yao}{Universite libre de Bruxelles}{A cohomogeneity-one hypersymplectic flow and some convergence results}

\begin{abspar}
Hypersymplectic flow is a dimensional reduction of the $G_2$ Laplacian flow to dimension four, which is designed to find Riemannian metrics with special holonomy groups. It has its own interest in the four dimensional symplectic geometry. In this talk, I will describe a cohomogeneity-one hypersymplectic flow on four-torus that exists and converges to the standard hyperK\"ahler structure. I will also discuss some convergence results for globally existing hypersymplectic flow under certain geometric assumptions. 

\end{abspar}

\abshead{Eiji Inoue}{University of Tokyo}{Moduli space of Fano manifolds admitting K\"ahler-Ricci solitons}

\begin{abspar}
Formulating a reasonable notion of moduli space of varieties 
is a fundamental problem in complex algebraic geometry. Moduli spaces of 
varieties are desired to be varieties (or complex spaces), as the notion 
of (possibly non-reduced) varieties are designed for describing 
deformations of themselves. 
In this talk, I will propose a reasonable notion of moduli space of Fano 
manifolds in view of K\"ahler-Ricci soliton, a kind of canonical metrics, 
and explain the key idea in the construction of the moduli space. We 
also quickly review some important results on K\"ahler-Ricci solitons, 
some of which are indispensable for our formulation. 

\end{abspar}

\abshead{Ngoc Cuong Nguyen}{POSTECH}{H\"older continuous subsolution theorem for the complex Monge-Amp\`ere equation}

\begin{abspar}
We solve the Dirichlet problem for the complex Monge-Amp\`ere equation on a strictly pseudoconvex domain with the right hand side being a positive Borel measure which is dominated by the Monge-Amp\`ere measure of a H\"older continuous plurisubharmonic function. If the boundary data is continuous, then the solution is continuous. If the boundary data is H\"older continuous, then so is the solution. In particular, the answer to a question of A. Zeriahi is always positive.

\end{abspar}

\dayhead{Sunday, July 15}

\abshead{Tien-Cuong Dinh}{National University of Singapore}{Sullivan's dictionary and ergodic theory for foliations}

\begin{abspar}
Sullivan's dictionary refers to the similarity between the theory of iteration of holomorphic maps, the dynamics of Kleinian groups and the theory of foliations by Riemann surfaces. In this talk, foliations by Riemann surfaces, possibly with singularities, will be considered. In particular, I will present recent results about hyperbolic entropy, ergodic theorem, heat equation and unique ergodicity. Some of the results do not confirm the expected analogies.
This talk is based on my works in collaboration with Viet-Anh Nguyen and Nessim Sibony.  

\end{abspar}

\clearpage


\abshead{Sean Curry}{UC San Diego}{Strictly pseudoconvex domains in $\mathbb{C}^2$ with obstruction flat boundary}

\begin{abspar}
On a bounded strictly pseudoconvex domain in $\mathbb{C}^n$, $n>1$, the smoothness of the Cheng-Yau solution to Fefferman's Monge-Amp\`ere equation up to the boundary is obstructed by a local curvature invariant of the boundary. For bounded domains in $\mathbb{C}^2$ diffeomorphic to the ball, we motivate and consider the problem of determining whether the global vanishing of this obstruction implies biholomorphic equivalence to the unit ball.

\end{abspar}

\abshead{Junwu Tu}{University of Missouri}{Gromov-Witten invariants of matrix factorizations}

\begin{abspar}
In this talk, I will describe a categorical construction of $g=0$ B-model Gromov-Witten invariants associated to a Landau-Ginzburg orbifold $(W,G)$.
One of the main ingredient of this construction is a categorical analogue of Saito's primitive forms. When the potential function $W$ is an invertible
polynomial and $G$ a diagonal symmetry group of $SL_n$, we prove that there exists a canonical choice of categorical primitive form for the category
$MF(W,G)$ of G-equivariant matrix factorizations of W. Conjecturally, these invariants are mirror to the $g=0$ part of the FJRW theory of the dual orbifold
$(W^t,G^t)$. This is joint work with Andrei Caldararu.

\end{abspar}

\abshead{Morimichi Kawasaki}{Kyoto University}{Poisson bracket invariant and the Lagrangian Floer theory}

\begin{abspar}
Polterovich introduced an Poisson bracket invariant defined for open
covers of symplectic manifolds. He gave its lower bound in some situations
by using the Hamiltonian Floer theory.
In this talk, I give another lower bound of his Poisson bracket invariant
in other situations by using the Lagrangian Floer theory. I also give its
application to the existence of non-displaceable fiber of moment maps.

\end{abspar}

%\clearpage

\dayhead{Monday, July 16}

\abshead{Yunhyung Cho}{Sungkyunkwan University}{Monotone Lagrangian tori in cotangent bundles of spheres}

\begin{abspar}
In this talk, we see how a sphere arises as a singular Lagrangian fiber of a Gelfand-Cetlin system of $U(n)$-type and $SO(n)$-type, respectively. Then, observing that $S^{2n+1}$ appears in Gelfand-Cetlin systems of both
types, we construct two families of Lagrangian tori in $T^{*}S^{2n+1}$ that are not pairwisely symplectomorphic
to each other, but have the same monotonicity constant. We also discuss their non-displaceablility by
computing the Floer cohomology. This is Joint work with Yoosik Kim and Yong-Geun Oh.

\end{abspar}
\clearpage

\abshead{Kei Irie}{University of Tokyo}{Chain level loop bracket and pseudo-holomorphic curves}

\begin{abspar}
Let $L$ be a Lagrangian submanifold in a symplectic vector space which 
is closed, oriented and spin. Using virtual fundamental chains of moduli 
spaces of nonconstant pseudo-holomorphic disks with boundaries on $L$, 
one can define a Maurer-Cartan element of a Lie bracket operation in 
string topology (the loop bracket) defined at chain level. This 
observation is due to Fukaya, who also pointed out its important 
consequences in symplectic topology. In this talk I will explain how to 
rigorously carry out this idea. Our argument is based on a string 
topology chain model previously introduced by the speaker, and the 
theory of Kuranishi structures on moduli spaces of pseudo-holomorphic 
disks, which has been developed by Fukaya-Oh-Ohta-Ono. 

\end{abspar}

%\clearpage

\abshead{Damian Brotbek}{University of Strasbourg}{Hyperbolicity properties of general complete intersections}

\begin{abspar}
A complex manifold is said to be Brody hyperbolic if it doesn't contain any entire curve.  Kobayashi conjectured in the 70's that if $H$ is a general hypersurface of sufficiently large degree in $\mathbb{P}^n$, then $H$ and its complement are both hyperbolic. 

One way to control the hyperbolicity properties of a variety is to study the positivity properties of its cotangent bundle. Indeed a variety with ample cotangent bundle is hyperbolic. It is known that the cotangent bundle of hypersurfaces in $\mathbb{P}^n$ is not ample, however, in a related direction, Debarre conjectured that that if $X$ is a general complete intersection of sufficiently large multi degree in $\mathbb{P}^n$ such that the codimension of $X$ exceeds its dimension, then the cotangent bundle of $X$ is ample. 

The relation between those two conjectured is best understood using jet differentials, higher order generalizations of symmetric differential forms.

In this talk we will present a method for constructing jet differential equations on some special complete intersection varieties and explain how this can be used to prove the conjectures of Debarre and Kobayashi.

\end{abspar}

%\clearpage

%\dayhead{Tuesday, July 18}

\abshead{Tatsuo Suwa}{Hokkaido University}{Representation of Sato hyperfunctions by relative Dolbeault classes}

\begin{abspar}
The \v{C}ech-de Rham cohomology is effectively used in various problems related to localization of characteristic classes. Likewise we may develop the \v{C}ech-Dolbeault cohomology theory and on the way we naturally come up with the relative Dolbeault cohomology. This cohomology may be interpreted as a notion dual to the mapping cone in the theory of derived categories and goes well with derived functors. In any case the cohomology turns out to be canonically isomorphic with the local (relative) cohomology of A. Grothendieck and M. Sato so that it provides a handy way of expressing the latter. 

In this talk I will present this cohomology theory and give, as applications, simple explicit expressions of the Sato hyperfunctions, some fundamental operations on them and related local duality theorems. The talk includes a joint work with N. Honda and T. Izawa. 

\end{abspar}
\clearpage
\abshead{Kyoung-Seog Lee}{IBS-CGP}{Special curves on some algebraic surfaces}

\begin{abspar}
Special curves on algebraic surfaces play important roles in arithmetic, birational and complex geometry. In this talk, I will discuss several special curves and their configurations on several explicitly given algebraic surfaces. Then I will discuss their roles in various contexts. Part of this talk is based on a joint work with JongHae Keum.
\end{abspar}

\dayhead{Tuesday, July 17}

\abshead{Hiroshi Iriyeh}{Ibaraki University}{Mahler's conjecture for three dimensional convex bodies and the rigidity of certain symplectic embeddings}


\begin{abspar}
Mahler's conjecture is one of the classical open problems in the area of convex geometry.
It states that for a centrally symmetric convex body K in the $n$-dimensional Euclidean space,
the product of the volume of K and that of the polar body is greater than or equal to $4^n/n!$.
The two dimensional case is solved by Mahler in 1939.
However, this conjecture is still open for $n>2$.
In this talk, we announce the solution of the conjecture for the three dimensional case.
We also discuss its application to Viterbo's isoperimetric-type conjecture for symplectic capacities
and a symplectic embedding problem.
The talk is based on a joint work with Masataka Shibata.

\end{abspar}

\abshead{Filippo Bracci}{University of Roma Tor Vergata}{Continuous extension of biholomorphisms in higher dimension}

\begin{abspar}
The boundary behavior of Riemann mappings in dimension one is completely understood in terms of the Carath\'eodory prime ends theory. In particular, every Riemann map whose image is a convex domain extends as a homeomorphism up to the boundary, except at most, two infinite singularities. Carath\'eodory's prime ends theory is related to quasi-conformality, and, with this condition, it holds in any dimension. However, biholomorphic maps in higher dimension are not in general quasi-conformal. In higher dimension, in fact, deep theories of extension of biholomorphisms (Fefferman's theorem, CR geometry) have been developed in case the domains have some regularity at the boundary. In a recent work with H. Gaussier, the speaker introduced a prime ends type theory, called “horosphere topology”, which allows to deal with non regular boundary. In particular, using this theory, one can prove that every convex map, with no boundary regularity assumed, from the ball (or from any bounded strongly pseudoconvex domain) extends as a homeomorphism except, at most, at two infinite singularities. The aim of this talk is to explain the previous result, and their generalizations, using Gromov's hyperbolicity theory and horosphere topology.

\end{abspar}



\end{document}
